% Folder 156-8, batch 03 — archivists' pages 41–60 (author's pp. 33–52).
% Pass: Opus 5.5 (claude-opus-5-5), 2026-10-02 — first pass, unchecked against the pages by a human.
%
% Notation fixed for this batch: |X| the support, |X|^\circ its "interior"
% (|X| minus the boundary), \tilde{X} the set of elements below X,
% \mathfrak{P}(P) the set of subsets, \Phi a figure (family of subsets),
% \mathcal{L} the ambient sets of the "Homomorphismes d'ateliers" section,
% \mathring{\ll} and \mathbin{|{\circ}|} his relations on multistructures.
\documentclass[11pt,a4paper]{article}
\input{../preamble/grothendieck}

\folder{156-8}
\batch{3}
\pages{41}{60}
\foldertitle{[Chapitre] VIII. Analysis situs (quatrième mouture) : notes manuscrites (26/06-04/07/1986).}
\dating{1986}
\watermark{Édition de démonstration}

\begin{document}

\section{Fonction support}

\page{41}\note{p.~33 de l'auteur ; sa pagination court sans interruption de 33 à 52 sur ce lot.}
J'en viens maintenant à la question d'une fonction genre « support »
\[
X \longmapsto |X| \qquad M \longrightarrow \mathfrak{P}(P),
\]
où $P$ est un ensemble (dont les éléments jouent le rôle de « points »). Voici
le minimum auquel je songe immédiatement (comparer \uncertain{CF}~VII,
p.~73)\note{lecture de « CF » douteuse ; peut-être « Ch » (chapitre). Même sigle p.~46.} :

\marginal{\uncertain{On veut} \ill{} $|A| = \bigcup_{X\in A} |X|$ \ill{} $M'$ \ill{}}

\begin{itemize}
  \item[(i)] $X \ll Y \Longrightarrow |X| \subset |Y|$ \marginal{concerne $\ll$}
  \item[\struck{(ii)}] $\big[$ (ii) $\forall X \in M$, $|X|^\circ \overset{\text{déf}}{=} |X| \smallsetminus |\partial X| \neq \emptyset$ $\big]$ \marginal{concerne $\leq$ \struck{(est impliqué par (iii))}}
  \item[(iii)] $\forall X, Y \in M$, $X \mathbin{|{\circ}|} Y \Longleftrightarrow |X|^\circ \cap |Y|^\circ = \emptyset$ \marginal{concerne $|{\circ}|$}
  \item[(iv)] (Pour \uncertain{normaliser} $X$) $|M| = P$, et les $|X|^\circ$ séparent les points de $X$, \ill{}
\end{itemize}
\note{la condition (ii) est enfermée par l'auteur entre crochets ; un premier début de ligne, \struck{(ii) $\forall X$ \ill{}}, est biffé.}

\textbf{NB} La condition (i) signifie : dire que $X \mapsto |X|$ est croissante
pour $\ll$, ou encore, pour $\mathring{\ll}$ et $\leq$. Si les $|X|$ séparent les
points de $X$\note{le $X$ est entouré.}, en est-il de même des $|X|^\circ$ ?
Soient $x \in |X|$, $y \notin |X|$. On voudrait trouver $Y \leq X$ avec
$x \in |Y|^\circ$, \uncertain{bien mieux encore} $y \notin |Y|^\circ$ puisque
$|Y|^\circ \subset |Y| \subset |X|$. Donc la question est de \struck{\ill{}}
savoir si $\exists\, Y \in \tilde{X}$ avec $x \in |Y|^\circ$. Donc on est amené à
renforcer (ii) en

\begin{itemize}
  \item[(ii bis)] $\forall X \in M$, la famille des $|Y|$, $Y \in \tilde{X}$, forme
  une figure \add{ens.}\ \struck{dans} (élémentaire) dans $P$, indexée par
  $\tilde{X}$, de support $|X|$.
\end{itemize}

\struck{Cela signifie \ill{} jusqu'à \ill{}} \add{Les sorts déjà
(\ill{} de (ii bis))} que $\tilde{X} \to \mathfrak{P}(P)$, $Y \mapsto |Y|$ est
injective, car si $Y, Y' \in \tilde{X}$ avec $Y \neq Y'$, on a
$Y \mathbin{|{\circ}|} Y'$ \add{(ii)} dès que $Y \neq Y'$, donc par (iii)
$|Y|^\circ \cap |Y'|^\circ = \emptyset$, donc par \struck{\ill{} on sait Je dis
que}\note{la phrase reste en suspens au bas de la page.}

\page{42}
A-t-on inversement, pour $Y, Z \in \tilde{X}$,
\[
|Y| \subset |Z| \Longrightarrow Y \leq Z \quad ??
\]
Ça ne semble pas conséquence de (i)(ii)(iii)(iv), \struck{À moins. Cela
L'est}. Est-il impliqué par (ii bis) ?? \struck{J'explicite} À vrai dire,
j'ai sous-entendu en y disant « indexé par $\tilde{X}$ » \ill{} comme un
\uncertain{ens. ordonné} (je présume) — mais peut-on le prouver ?
\struck{Soit $L = \tilde{Y} \cap \tilde{Z}$} Ce qui est sûr, c'est que j'aurai
besoin de savoir que, pour $Y \in \tilde{X}$, $|Y|^\circ$ est aussi l'ens.
\struck{des} \ill{} que, défini en termes de la figure ens. indexée par
$\tilde{X}$, i.e. que
\[
|Y|^\circ \overset{?}{=} |Y| \smallsetminus \bigcup_{\substack{Z \in \tilde{X} \text{ tel que}\\ |Z| \subsetneq |Y|}} |Z| \overset{\text{déf}}{=} |Y|^{\circ\prime}
\]
mais dans la $\bigcup$, il y a a priori \add{peut-être} plus des $Z$ que
\struck{ceux} les seuls $Z < Y$, donc le \ill{} risque d'être plus petit
\add{(\uncertain{forcément égal ?})} que le \struck{précédent}. Mais (ii bis)
(\ill{} définition) \ill{} que les $|Y|^{\circ\prime}$ recouvrent déjà $|X|$ (et
sont non vides). A fortiori les $|Y|^\circ$ recouvrent \struck{donc}. Mais on
sait \add{qu'ils sont deux à deux} disjoints par (iii) — d'où finalement
l'égalité $|Y|^\circ = |Y|^{\circ\prime}$, et la relation
\[
|X| = \bigcup_{Y \in \tilde{X}} |Y|^\circ .
\]
Revenons donc à la question si
\[
|Y| \subset |Z| \Longrightarrow Y \leq Z ,
\]
introduisons
$\tilde{Y} \cap \tilde{Z} = L$, \struck{\ill{}} $\tilde{Y} \cup \tilde{Z} = M$.

\page{43}
\[
\text{\struck{$|Y| = \coprod_{Z \in \tilde{Y}} |Z|^\circ = \bigcup_{Z \in L}$}}
\]
Comme les $|Z|^\circ$ ($Z \in M$) sont \add{non vides et} deux à deux
disjoints, \struck{et \ill{}} et on a\note{dans les réunions qui suivent,
l'indice muet est écrit avec un $Z$ barré, que nous rendons par $Z$ ; le
$|Z|$ de gauche de la deuxième ligne est le $Z$ fixé.}
\marginal{ici on utilise (ii bis) non seulement pour $X$, mais aussi pour $Y$ et $Z$}
\[
|Y| = \bigcup_{Z \in \tilde{Y}} |Z|^\circ = \Big(\bigcup_{Z \in L} |Z|^\circ\Big) \cup \Big(\bigcup_{Z \in \tilde{Y} \smallsetminus L} |Z|^\circ\Big)
\]
\struck{$|Z| =$} et de même
\[
|Z| = \Big(\bigcup_{Z \in L} |Z|^\circ\Big) \cup \Big(\bigcup_{Z \in \tilde{Z} \smallsetminus L} |T|^\circ\Big)
\]
\note{le dernier terme porte bien $|T|^\circ$ sur la page.}
d'où
\[
|Y| \cap |Z| = \bigcup_{Z \in L} |Z|^\circ ,
\]
et on \uncertain{voit} \uncertain{égal à} $|Y|$ \ill{} i.e. $|Y| \subset |Z|$,
que si la \struck{\ill{}} \ill{} $\bigcup_{Z \in \tilde{Y} \smallsetminus L} |Z|^\circ \neq \emptyset$
\note{on attendrait « $= \emptyset$ » ; nous lisons « $\neq$ ».}
i.e. $\tilde{Y} \smallsetminus L = \emptyset$, i.e. $Y \leq Z$,
\struck{\ill{}} on gagne ! Ainsi, on a \add{(pour $X$ et les $Y \in \tilde{X}$)}
\ill{} que (ii bis) \add{(implique)}

\begin{itemize}
  \item[(ii ter)] $\displaystyle |X| = \bigcup_{Y \in \tilde{X}} |Y|^\circ$
  \note{sous la réunion l'auteur a d'abord écrit $|X|^\circ$, puis entouré le $X$ et l'a corrigé en $Y$ par une flèche.}
\end{itemize}

(\struck{\ill{}} c'est une partition \add{indexée par $\tilde{X}$}, par (ii) et
(iii), impliquant que $Y \neq Z$ ($Y, Z \in \tilde{X}$) $\Rightarrow
|Y|^\circ \cap |Z|^\circ = \emptyset$ donc $|Y|^\circ \neq |Z|^\circ$
(puisque $|Y|^\circ \neq \emptyset$ par (ii)))

Inversement, si on suppose (ii ter), a-t-on (ii bis) ? \struck{Voyons
\uncertain{CF}~V p.~1} \struck{Prop.~1 condition 3° ($\alpha$, $\beta$,
$\gamma$). La \ill{}} \ill{} $\gamma$ \struck{déjà vu la \ill{} \ill{} ici,
et en celui déduit de la} \ill{} avec (i) les $|X|^\circ \neq \emptyset$
(\textbf{NB} l'identité des $|X|^\circ$ défini ici, et de celui défini par la
famille des $|Y|$ ($Y \in \tilde{X}$) est vraie \ill{} \add{(cf deux lignes plus
haut) \ill{}\,…}) et (ii) $\forall x \in |X|$, $\exists$ un plus petit
$Y \in \tilde{X}$ avec $x \in |Y|^\circ$. Que la relation d'ordre de
$\tilde{X}$ soit celle induite par l'inclusion dans $\mathfrak{P}(P)$, on l'a
vu par l'argument
\note{passage très raturé : un premier énoncé biffé, puis un « NB » encadré
entre deux traits ; la lecture ci-dessus en donne la suite des idées, non
l'ordre exact des ajouts.}

\page{44}
précédent, qui a utilisé seulement le fait que les $|Y|$ sont les réunions
disjointes des $|Z|^\circ$ pour $Z \in \tilde{Y}$.) \struck{Soit donc $x$.
Comme Il faut \ill{}} Soit donc $Y \in \tilde{X}$ avec $x \in |Y|^\circ$.
Prouvons que $Y$ est le plus petit él. de $\tilde{X}$ tel que $x \in |Z|$,
i.e.
\[
x \in |Y|^\circ,\ x \in |Z| \Longrightarrow Y \leq Z .
\]
Cela se voit \uncertain{toujours} par le même argument que tantôt :
\uncertain{puisque} donc ici $Y \in \tilde{Z}$ donc $Y \leq Z$.

Dans \struck{tous} ces cas, la relation $\ll$ n'a \uncertain{pas joué de}
rôle, \ill{} on a travaillé dans l'ens. ordonné \ill{}. Donc il faut
expliciter la \ill{}

\note{de « Donc il faut » jusqu'au bas de la page, le texte est barré d'un
long trait oblique ; nous le donnons tel qu'il se lit sous la rature.}

\struck{\textbf{Lemme} Soit $I$ un ens. ordonné, $P$ un ens.,}
\[
\text{\struck{$\varphi : I \longrightarrow \mathfrak{P}(P), \quad i \longmapsto X_i \text{ ou } |i|$}}
\]
\struck{une application, \struck{qui soit \ill{}} satisfaisant les conditions
(on a posé $X_i^\circ = X_i \smallsetminus \bigcup_{j<i} X_j$)}
\begin{itemize}
  \item[(i)] \struck{croissante : $i \leq j \Rightarrow X_i \subset X_j$}
  \item[(ii)] \struck{$\forall i, j$, $i \neq j \Longrightarrow X_i^\circ \cap X_j^\circ = \emptyset$, puis on demande que $\forall i$ $X_i^\circ \neq \emptyset$}
\end{itemize}
\note{(ii) a d'abord été écrit « $\forall i,j \in I$, $i \neq j \Longleftrightarrow |X_i|^\circ \cap |X_j|^\circ = \emptyset$ », puis corrigé en surcharge.}
\struck{(Alors $\varphi$ \add{est injective}). Alors les conditions suivantes
sont équivalentes}
\begin{itemize}
  \item[(iii)] \struck{\struck{L'image de} $\varphi(I) = I' \subset \mathfrak{P}(P)$ est une figure ensembliste dans $P$, et $\forall i$ il en est de même \add{de l'image} de $I_{\leq i}$ ou $I_{<i}$}
  \item[(iii bis)] \struck{$\forall i$, $X_i = \bigcup_{j \leq i} X_j^\circ$}
\end{itemize}
\struck{Les conditions \ill{} impliquent que $\varphi$ \add{induit} un isom.
d'ens. ordonnés $I \xrightarrow{\sim} \tilde{I}$}

\section{Figures indexées}

\page{45}
Je vais reprendre ici la prop.~1 p.~1 de \uncertain{CF}~V, dans un contexte
« indiciel » qui sera, en fait, un peu plus général.

\textbf{Proposition} Soient $I$ un ens. ordonné, \struck{\ill{}} $X$ un
ensemble \uncertain{(c.-à-d.} \ill{} ce qui désignait $P$ dans
les pages précédentes), $\Phi = (X_i)_{i \in I}$ une famille d'ens. indexée
par $I$. On pose, $\forall i \in I$
\marginal{On suppose $i \leq j \Rightarrow X_i \subset X_j$ [je ne suppose \ill{} \ill{}, mais \ill{} pour définir $X_i^\circ = X_i \smallsetminus \partial X_i$ \ill{} \ill{} $\partial X_i \subset X_i$)}
\begin{align*}
\partial X_i &= \bigcup_{j<i} X_j \;\subset X_i \\
X_i^\circ &= X_i \smallsetminus \partial X_i \\
|\Phi| &= \bigcup_{i \in I} X_i .
\end{align*}
On note que $\forall x \in X$, \struck{$i \in I$}, si on pose
\[
I^x = \{ i \in I \mid x \in X_i \}
\]
on a l'équivalence
\[
x \in X_i^\circ \Longleftrightarrow i \text{ est un élément minimal de } I^x .
\]
Conditions équivalentes
\begin{itemize}
  \item[(i)] $\forall x \in |\Phi|$, l'ensemble $I^x$ a un plus petit élément.
  \item[(ii)] \struck{Pour} tout $x \in |\Phi|$, l'ensemble $I^x$ a un élément
  minimal, \add{(\uncertain{automat.} si $I$ fini)} et de plus pour $i, j \in I$, on a
  \[
  X_i \cap X_j = \bigcup_{k \leq i, j} X_k .
  \]
  \item[(iii)] On a les conditions suivantes
\end{itemize}
\begin{itemize}
  \item[a)] $\forall i, j \in I$ avec $i \neq j$, on a $X_i^\circ \cap X_j^\circ = \emptyset$
  \item[b)] \struck{\ill{}} $\forall i \in I$, on a
  \[
  X_i = \bigcup_{j \leq i} X_j^\circ \qquad \text{(automat. si $I$ fini)}
  \]
\end{itemize}
\begin{itemize}
  \item[(iv)] Posant \struck{(cas \textbf{NB} a) et b) impliquent que la famille
  \ill{} pour}
  \[
  I' = \{ i \in I \mid X_i^\circ \neq \emptyset \} ,
  \]
  on a : \struck{alors la famille des}
\end{itemize}
\begin{itemize}
  \item[$\alpha$)] La famille $\{X_i^\circ\}_{i \in I'}$ forme une
  partition de $|\Phi|$, \struck{\ill{}}
\end{itemize}
\note{la condition (iv) se poursuit p.~46.}

\page{46}
\marginal{\textbf{Question :} $\beta$) et $\gamma$) conséquences de $\alpha$) ? Oui si $I$ fini, cf p.~44, mais \uncertain{exemple} \ill{} $\alpha$) et $\gamma$) \ill{}, mais \ill{} $\beta$) ? Non, $\beta$) \uncertain{conséquence} \ill{} \ill{} \ill{}}
\begin{itemize}
  \item[$\beta$)] Pour tout $i \in I$, $X_i$ est saturé par la relation
  \struck{d'équivalence} d'équivalence définie par cette partition.
  \item[$\gamma$)] Pour $i, j \in I$, \struck{avec $X_j^\circ \neq 0$, la
  relation} \add{et $X_j^\circ \neq \emptyset$} on a
  \[
  X_j^\circ \subset X_i \Longrightarrow j \leq i .
  \]
\end{itemize}
\note{dans $\gamma$) l'indice de droite a été surchargé ; nous lisons $i$.}

\textbf{Disons} \struck{la} condition (i), \struck{i.e. \ill{} des $I^x$
signifie} \add{pour tout $x \in |\Phi|$, $I^x$} \struck{(qu'il} ait un plus
petit élément), signifie ($\forall x \in X$) trois choses \struck{trois
choses} \add{(dont on sait qu'il est $\neq \emptyset$)} :
\marginal{évidemment $3°) \Rightarrow 1°)$ (car on sait $I^x \neq \emptyset$)}
\begin{itemize}
  \item[1°)] $I^x$ a un élément minimal, $i$, i.e. $\exists\, i \in I$ tel que
  \struck{\ill{}} $i$ est élt. minim. de $I^x$ i.e.
  \[
  x \in X_i^\circ .
  \]
  \item[2°)] Deux él. minimaux de $I^x$ sont égaux i.e.
  \[
  \forall i, j \in I,\ x \in X_i^\circ \cap X_j^\circ \Longrightarrow i = j
  \]
  \item[3°)] Tout élément de $I^x$ est minoré par un élément minimal,
  \struck{\ill{}} i.e.
  \[
  \forall i \in I \text{ tel que } x \in X_i,\ \exists\, j \leq i \text{ tel que } x \in X_j^\circ .
  \]
\end{itemize}
C'est donc la conjonction des trois conditions : $\forall x$ on a 1°,
$\forall x$ on a 2°, $\forall x$ on a 3°. Cela s'écrivent
\begin{itemize}
  \item[1')] $|\Phi| = \bigcup X_i^\circ$ \quad (automat. si $I$ fini)
  \item[2')] $\big[X_i^\circ \cap X_j^\circ \neq \emptyset \Rightarrow i = j$ i.e.$\big]$ $i \neq j \Longrightarrow X_i^\circ \cap X_j^\circ = \emptyset$
  \item[3')] $\forall i \in I$, on a $X_i = \bigcup_{j \leq i} X_j^\circ$ \quad (automat. si $I$ fini)
\end{itemize}
On voit d'ailleurs, puisque $|\Phi| = \bigcup_{i \in I} X_i$, que 3') implique
1'). Donc la condition (i) équivaut à la conjonction de 2' et 3'), qui ne sont
autres que a), b) de (iii). Ainsi
\[
\text{(i)} \Longleftrightarrow \text{(iii)} .
\]
\marginal{\textbf{NB} 3', 3' \ill{} $I^x$ \ill{} \ill{} \ill{} $X_i^\circ \cap X_j^\circ$ \ill{} — \textbf{Question :} si $\forall i,j$, $i \neq j$ \ill{} ? Oui si $I$ fini, Non, cf p.~44}
D'autre part, les parties $\alpha$), $\beta$), $\gamma$) de (iv) se
traduisent respectivement en :
\begin{itemize}
  \item[1'')] $|\Phi| = \bigcup_{i \in I'} X_i^\circ$ \quad (équivalent à 1') \struck{et $i, j \in I'$}
  \item[\struck{2'')}] \struck{$\forall i \in I$, posant $I(i) =$}
\end{itemize}
et $\forall i, j \in I'$, \struck{$X_i^\circ \cap X_j$} on a
$X_i^\circ \cap X_j^\circ = \emptyset$ (équivalent à 2').

\page{47}
\begin{itemize}
  \item[2'')] $\forall i \in I$, posant
  \[
  I'(i) = \{ j \in I' \mid X_j^\circ \subset X_i \} ,
  \]
  \[
  \text{\struck{$I(i) = \{ j \in I \mid X_j^\circ \subset X_i \}$}}
  \]
  \struck{(on note que $I'(i) = I(i) \cap I'$, et $\bigcup_{j \in I(i)} X_j^\circ = \bigcup_{j \in I'(i)} X_j^\circ$)},
  on a
  \[
  X_i = \bigcup_{j \in I'(i)} X_j^\circ
  \]
  \struck{on a qui \uncertain{revient au} même $X_i = \bigcup_{j \in I'(i)} X_j^\circ$, i.e.~3')}
  \item[3'')] On a, pour tout $i \in I$
  \[
  I'(i) = I' \cap I_{\leq i}
  \]
  \marginal{ce qui implique que $\bigcup_{j \in I'(i)} X_j^\circ = \bigcup_{j \leq i} X_j^\circ$}
\end{itemize}
On voit donc que moyennant 3''), 2'') équivaut à \struck{1'', 2'', 3'',} 3').
Ceci dit, l'ens. des conditions 1'', 2'', 3'' équivaut aussi à celui-ci
\[
\left\{
\begin{array}{lll}
|\Phi| = \bigcup_{i \in I} X_i^\circ & & \text{i.e. (1')} \\
\forall i, j \in I & X_i^\circ \cap X_j^\circ = \emptyset & \text{i.e. (2')} \\
\forall i \in I & X_i = \bigcup_{j \leq i} X_j^\circ & \text{i.e. (3')} \\
\end{array}
\right.
\]
$+$ la condition $\gamma$) \add{et $X_j^\circ \neq \emptyset$} :
\[
\forall i, j \in I \quad X_j^\circ \subset X_i \Longrightarrow j \leq i .
\qquad 3'') \Leftrightarrow \gamma)
\]
\note{dans la deuxième ligne l'auteur omet ici l'hypothèse $i \neq j$.}
Cela montre donc que
\[
\text{(iv)} \Longleftrightarrow \text{(i)} + \gamma)
\]
et l'équivalence (iv) $\Longleftrightarrow$ (i) revient \struck{donc} à
(i) $\Longrightarrow$ ($\gamma$) qui est contenu dans le

\marginal{Attention ! \ill{} $X_j \subset X_i$ \ill{} conditions (avec $X_j^\circ \neq \emptyset$) \ill{} $X_i \supset X_j^\circ$ (\ill{} \ill{} si (i) est \ill{}) \ill{} $\Phi$ \ill{} !! \ill{} \ill{} \ill{}}

\textbf{Cor.\ \uncertain{Main}~1} Sous les conditions équivalentes (i), (iii),
\struck{\ill{} tel que $X_j^\circ \neq 0$ les} \add{\struck{(équivalentes} puis
\ill{}} on a, pour $i, j \in I$, \struck{Conditions \ill{}}
\[
(j \leq i) \Longleftrightarrow (X_j \subset X_i) \Longleftrightarrow (X_j^\circ \subset X_i) \Longleftrightarrow (X_j^\circ \cap X_i \neq \emptyset)
\]
\note{les indices de cette ligne sont surchargés à plusieurs reprises ; une
variante biffée au-dessus de la dernière équivalence est illisible.}

(\textbf{NB} Dans l'énoncé de la prop., on n'a pas supposé que
$i \leq j \Rightarrow X_i \subset X_j$ ?!)

\textbf{Dém. du Cor.~1} en vertu de 3') on a
\[
X_i = \bigcup_{i' \leq i} X_{i'}^\circ , \qquad X_j = \bigcup_{j' \leq j} X_{j'}^\circ
\]
donc
\[
i \leq j \Longrightarrow X_i \subset X_j \Longrightarrow \text{\struck{\ill{}}}
\quad \text{(si on suppose $X_j^\circ \neq \emptyset$)}
\]
\note{la ligne se poursuit p.~48.}

\page{48}
D'ailleurs les implications
\[
(X_j \subset X_i) \Longrightarrow (X_j^\circ \subset X_i) \overset{\text{si } X_j^\circ \neq \emptyset}{\Longrightarrow} (X_j^\circ \cap X_i)
\]
sont triviales. Reste il reste à prouver, si $X_j^\circ \neq \emptyset$
\[
\text{\struck{$X_j^\circ \cap X_i \Longrightarrow i \leq j$}}
\]
\[
X_j^\circ \cap X_i \neq \emptyset \Longrightarrow j \leq i .
\]
\note{au-dessus de $X_i$ dans la ligne biffée, l'auteur a écrit $X_i^\circ$.}
Mais $X_i = \bigcup_{i' \in I'(i)} X_{i'}^\circ$ \struck{est saturé pour}
\add{par 3')}, et comme les $X_{i'}^\circ$ ($i' \in I'$) forment une partition
de $|\Phi|$, \struck{\ill{}} \add{la \uncertain{partition}}
$X_j^\circ$ (pour $j \in I'$ donc) ne peut rencontrer $X_i$ que si
$j \in$ \struck{$I'(i)$} $I'(i)$ i.e.\ $j \leq i$, cqfd.

Ainsi on a prouvé
\[
\text{(i)} \Longleftrightarrow \text{(iii)} \Longleftrightarrow \text{(iv)}
\]
on a même tenu (?) le cor.~Main, dont je signale tout de suite le
(dans les conditions équiv.\ de la prop.)

\textbf{Cor.\ \uncertain{Main}~2} Supposons \add{(que $\forall i \in I$, on ait $X_i^\circ \neq \emptyset$.)}
Alors l'application
\[
i \longmapsto X_i : \varphi : I \longrightarrow \mathfrak{P}(X)
\]
est injective, et \struck{est} \add{est un plongement} d'ensembles ordonnés :
$\forall i, j \in I$, on a
\[
X_i \subset X_j \Longleftrightarrow i \leq j .
\]
Ainsi, les conditions envisagées (i) \struck{\ill{}} \add{plus $\forall i \in I$,
$X_i^\circ \neq \emptyset$,} \add{sont} \struck{\ill{}} nécessaires et suff.\
pour que $\varphi$ \struck{\ill{}} \add{soit} \struck{une} $(X_i)_{i \in I}$
\struck{\ill{} soit} une figure ensembliste indexée par l'ens.\ ord.\ $I$.
\struck{(\ill{} \ill{} !?)}

Il reste à voir l'équivalence avec (ii).

\struck{La condition (iii b) \ill{} que} \struck{Le fait que \ill{} $I^x$
\add{(} ait un élt. minimal \add{($x \in |\Phi|$)} équivaut à 3')
revient à l'\ill{} : $|\Phi| = \bigcup_{i \in I} X_i^\circ$, i.e.~1').}
\note{ces dernières lignes sont barrées de plusieurs traits obliques.}

\page{49}
(i) $\Longrightarrow$ (ii). Il reste à prouver la formule
\[
X_i \cap X_j = \bigcup_{k \leq i, j} X_k
\]
(où $\supset$ est évidente). Mais par 3') on a
\[
X_i = \bigcup_{i' \leq i} X_{i'}^\circ , \qquad X_j = \bigcup_{j' \leq j} X_{j'}^\circ
\]
d'où
\[
X_i \cap X_j = \bigcup_{\substack{i' \leq i \\ j' \leq j}} X_{i'}^\circ \cap X_{j'}^\circ .
\]
Or par 2') on a $X_{i'}^\circ \cap X_{j'}^\circ = \emptyset$ \uncertain{toujours}
si $i' \neq j'$, d'où
\[
X_i \cap X_j = \bigcup_{k \leq i, j} X_k^\circ \subset \bigcup_{k \leq i, j} X_k
\qquad \text{cqfd}
\]

(ii) $\Longrightarrow$ (i). Cela résulte du

\textbf{Lemme} \struck{Conditions} Supposons qu'on ait, $\forall i, j \in I$
\[
X_i \cap X_j = \bigcup_{k \leq i, j} X_k .
\]
Alors on a les propriétés :
\marginal{Si de plus $\bigcup_{i \in I} X_i^\circ = |\Phi|$ on a \ill{}}
\begin{itemize}
  \item[2')] $\forall i, j \in I$, $i \neq j$ \struck{\ill{}} :\quad $X_i^\circ \cap X_j^\circ = \emptyset$
  \item[\struck{3')}] \struck{$\forall i \in I$ :\quad $X_i = \bigcup_{j \leq i} X_j^\circ$}
\end{itemize}
\note{la condition 3') est barrée d'un trait qui coupe aussi la fin de 2') ; seul 2') est démontré ensuite.}

Prouvons 2') : on a
\[
X_i^\circ \cap X_j^\circ \subset X_i \cap X_j = \bigcup_{k \leq i, j} X_k
\]
si $i \neq j$, \struck{et} et les éléments $k$, tous $k \leq i, j$ \add{sont} ou
bien $< i$, ou bien $< j$ (vu que \uncertain{les deux} égalités $k = i$ et $k = j$
s'excluent) donc
($k$ \ill{} $X_k \subset \partial X_i$ ou $X_k \subset \partial X_j$, donc
\[
i \neq j \Longrightarrow X_i \cap X_j \subset \partial X_i \cup \partial X_j \subset \complement (X_i^\circ \cap X_j^\circ)
\]
donc $X_i^\circ \cap X_j^\circ \subset \complement X_i^\circ \cap X_j^\circ$
d'où $X_i^\circ \cap X_j^\circ = \emptyset$ \add{(ou bien \uncertain{moins} $\supset$, il faut aussi $\subset$.)}

Prouvons 3'). (Soit $x \in X_i$, par hyp.\ $\exists\, j \in I$ tel que
$x \in X_j^\circ$, tout revient à voir que $X_j^\circ \subset X_i$
\note{la phrase se poursuit p.~50.}

\page{50}
i.e.\ la susdite implication
\[
X_j^\circ \cap X_i \neq \emptyset \Longrightarrow j \leq i .
\]
Mais on a
\[
X_j^\circ \cap X_i \subset X_j \cap X_i = \bigcup_{k \leq i, j} X_k
\]
et si on n'avait pas $j \leq i$, on ne \uncertain{trouverait} dans la
réunion des derniers membres aucun $k = j$, donc cette réunion est
$\subset \bigcup_{k < j} X_k = \partial X_j \subset \complement X_j^\circ$ \add{$\complement (X_j^\circ \cap X_i)$ donc} :
$X_j^\circ \cap X_i$ est contenu dans son complémentaire, ce qui
\uncertain{contredit} l'hyp.\ que c'est vide.
\note{lire sans doute « non vide » ; l'hypothèse est $X_j^\circ \cap X_i \neq \emptyset$.}

La prop.\ est prouvée.

\textbf{Cor.\ \uncertain{Main}~3} Pour que $(X_i)_{i \in I}$ soit une figure
ensembliste indexée par l'ens.\ ordonné $I$ (cf Cor.~2), il faut et il suffit
qu'on ait les conditions suivantes
\begin{itemize}
  \item[a)] $\forall i, j \in I$ \struck{avec $i \leq j$, on a $X_i \subset X_j$} \struck{$X_i^\circ \neq \emptyset$, $i \neq j \Leftrightarrow$ \ill{} $i \neq j$} \add{$X_i^\circ \cap X_j^\circ = \emptyset$}
  \item[b)] $\forall i, j \in I$, on a \ill{} \ill{} $\Phi_i \subset \mathfrak{P}(X)$
  \item[c)] $\forall i \in I$, si on désigne par $\Phi_i$ \ldots\ l'ens.\ des parties de $X$ de la forme $X_j$, $j \leq i$, $\Phi_i$ est une figure ensembliste.
\end{itemize}
\note{l'ordre des lignes b) et c) est incertain : les lettres sont surchargées et la ligne b) semble se fondre dans c).}
\marginal{\textbf{NB} b) équivaut à \ill{} \ill{} $X_i^\circ \neq \emptyset$, $X_i^\circ \cap X_j^\circ = \emptyset$ \ldots\ b') $\forall i \in I$, b'') $\forall i, j$, $i \neq j$)}
\marginal{\textbf{NB} La condition b) \ill{} implique déjà $\bigcup X_i^\circ$ \ill{}}

\struck{On trouve} $\Phi_i$ est une figure \uncertain{\og quasi-ordonnée \fg},
et moins encore, \struck{(Il suffit = qu'elle soit} \struck{(\ill{} \ill{} $\forall i$}
que \struck{tout} $x \in X_i$, \add{dans l'ens.\ des $X_j$ ($j \leq i$) qui le
contiennent, il y ait un élément minimal} \ill{} qu'il $X_i$ \ill{} i.e.\
$I$ \ill{} $F$ !!) \ill{} conditions \ill{} ; prouvons $3'$) que
$X_i = \bigcup_{j \leq i} X_j^\circ$ \struck{(\ill{} \ill{} \ill{} 2'), \ill{} 2') $= b''$)}.

\struck{a) b) implique \ill{} $\forall i, j \in I$, \ill{}}
\struck{$X_i \subset X_j \Longleftrightarrow i \leq j$ ($\Leftarrow$ est contenu par a))}

\struck{Ce qui implique 2') et 3') donc (i).}

\struck{Donc, $\forall i \in I$, $X_i^\circ =$ \uncertain{intérieur} \ill{} de
$X_i$ dans $\Phi_i$ $= X_i \smallsetminus \bigcup_{j \leq i,\ X_j \neq X_i} X_j$}
\note{le bas de la page est barré de longs traits obliques.}

\page{51}
Fixons un $i \in I$. Pour tout $j \leq i$, posons
\[
X_j^{\circ\prime} = \text{cellule ouverte de } X_j \text{ dans } \Phi_i
= X_j \smallsetminus \bigcup_{\substack{j' \leq i \\ \text{t.q. } X_{j'} \subsetneq X_j}} X_{j'}
\]
\marginal{\textbf{NB} dans la somme, on ne suppose pas $j' \leq j$}
\add{Je dis que} cette somme comporte plus de termes que la somme analogue
qui figure dans la déf.\ de
\[
X_j^\circ = X_j \smallsetminus \bigcup_{j' < j} X_{j'} ,
\]
i.e.\ que
\[
j' < j \Longrightarrow X_{j'} \subsetneq X_j .
\]
En effet, on \struck{\ill{} $X_j$} $X_{j'} \subset \partial X_j$, donc
$X_j \smallsetminus X_{j'} \supset X_j \smallsetminus \partial X_j = X_j^\circ$,
comme \struck{\ill{}} donc on conclut que bien $X_j \smallsetminus X_{j'} \neq \emptyset$
i.e.\ $X_j^\circ \neq \emptyset$ par hyp.\ on conclut $X_{j'} \neq X_j$.

On a donc prouvé par là-même
\[
X_j^{\circ\prime} \subset X_j^\circ
\]
\marginal{donc $\bigcup X_j^{\circ\prime} = X_i$ \ill{} $\bigcup X_j^\circ = X_i$, \uncertain{cqfd}.}
(«~figure~») \add{(par la \uncertain{cellule} :)}. \struck{Or} \struck{Mais} Mais par hyp.\ \ill{}
l'une (des $X_j^{\circ\prime}$) recouvre \add{\uncertain{tout}} \struck{donc}
\[
\bigcup_{j \leq i} X_j^{\circ\prime} \subset \bigcup_{j \leq i} X_j^\circ \subset X_i
\]
avec égalité des deux extrémités. Comme les $X_j^\circ$ sont $\neq \emptyset$ et
deux à deux disjoints, cela implique que
\[
X_j^{\circ\prime} = X_j^\circ :
\]
\struck{les cellules ouvertes sont les $X_j^\circ$.} \add{En même}
\struck{En même temps,} \struck{\ill{} \ill{} que} \add{temps,} si $j, k \leq i$,
je dis que \struck{$j \neq k \Rightarrow$}

\struck{Soient $i, j \in I$ tels que $X_i \subset X_j$, prouvons $i \leq j$. Soit
$x \in X_i^\circ$ ($X_i^\circ \neq \emptyset$ par b'))}
\note{toute la seconde moitié de la page, depuis « On a donc prouvé », est
barrée de longs traits obliques ; la fin, en particulier, est entourée et
biffée.}

\page{52}
\struck{par \ill{} $\exists\, j' \leq j$ avec $x \in X_{j'}^{\circ\prime}$, a fortiori $x \in X_{j'}^\circ$.
Donc $X_i^\circ \cap X_{j'}^\circ \neq \emptyset$, \uncertain{d'où} par b''), on a $i = j'$,
d'où $i \leq j$, cqfd.}
\note{ces trois lignes, suite de la page précédente, sont barrées de traits obliques.}

\textbf{Exemple} Voici un exemple où on a 1' et 2', et en plus tous les
$X_i^\circ \neq \emptyset$, sans avoir 3') (\uncertain{indexations} respectant la
relation d'inclusion $i \leq j \Rightarrow X_i \subset X_j$).
Soit \struck{$X =$} $A_0 \supset A_1 \supset \cdots \supset A_n \supset \cdots$
une suite strictement décroissante d'ens.\ dans $X$, avec
\[
\text{\struck{$A_\infty$}}\ N = \bigcap A_n \neq \emptyset .
\]
\struck{$B = A_n$}
Soit $A_\infty \subset X$ avec $A_\infty \cap A_0 = N$, $A_\infty \not\subset A_0$
i.e.\ $N \neq A_\infty$. On prend $I \subset \mathfrak{P}(X)$ défini comme
\add{$I \smallsetminus \{\infty\}$} \struck{l'ens.} \ill{} \ill{} ordonné par
inclusion, mais $\infty$ est un «~pt isolé~». On a
\[
A_i^\circ = A_i \smallsetminus A_{i+1} \neq \emptyset \quad \text{si } i \neq \infty
\]
\[
A_\infty^\circ = A_\infty \neq \emptyset .
\]
\struck{\ill{}} Les $A_i^\circ$ ($i \in I$) sont mutuellement disjoints non
vides, leur réunion est $A_0 \cup A_\infty$ ($= |\Phi|$). Mais un $A_i$
($i \neq \infty$) n'est pas réunion des $A_j^\circ$ avec $A_j \subset A_i$
(il manque $N$ !), il n'est \uncertain{pas} saturé par la relation d'équivalence
définie par les $A_j^\circ$ ($j \in I$). \struck{D'autre part,} \add{Un}
$x \in N$ n'est pas contenu dans un plus petit élément de $\Phi'$
\add{($= N$)}. On a $A_\infty^\circ \cap A_i \neq \emptyset$
\add{$\overset{\shortparallel}{A_\infty \cap A_i}$}, sans pour autant avoir
$A_\infty^\circ \subset A_i$ \struck{\ill{} bien $A_\infty^\circ \subset A_i$}.

Ici la condition (iv $\beta$) n'est pas satisfaite (\ill{} des $A_i$), alors
que $\alpha$, $\gamma$) le sont.

\textbf{Question :} si on suppose,
\[
\left\{
\begin{array}{l}
\forall i \quad X_i^\circ \neq \emptyset \\
\forall i \neq j \quad X_i^\circ \cap X_j^\circ = \emptyset \\
\text{\struck{$\bigcup X_i^\circ = |\Phi|$}}
\end{array}
\right.
\]
$+$ la saturation des \uncertain{dits} $X_i$,
\[
X_i = \bigcup_{j \text{ t.q. } X_j^\circ \subset X_i} X_j
\]

\page{53}
a-t-on $\gamma$) i.e.\ savoir que
\[
X_j^\circ \subset X_i \Longrightarrow j \leq i
\]
dans tous ? Du moins si \struck{l'indexation} on suppose
\[
X_j \subset X_i \Longrightarrow j \leq i
\]
(«~indexations respectant les inclusions~»)
i.e.
\[
X_j^\circ \subset X_i \Longrightarrow \text{\struck{$j \leq$}}\ X_j \subset X_i \ ?
\]
Voici un autre contre-exemple, \uncertain{moins} \uncertain{idiot}.
\struck{Supposons qu'on ait deux familles str.\ décroissantes, $(A_i)$ comme
dessus et $B_0 \supset B_1 \supset \cdots B_n \supset \cdots$ Supposons
$A_0 \cap B_0 = \bigcap A_i = \bigcap B_i$, soit $N$ et prenons
$\Phi = \{ A_i, B_j \}$. On} \add{$A_i$}
\note{ce premier essai est barré de longs traits obliques ; l'auteur reprend
avec la seule famille $(A_i)$.}

Reprenons\add{\uncertain{\,:}} la famille un dernier ensemble
\[
A = A_\infty - N
\]
i.e.\ prenons
\[
\Phi = \{ (A_i), A_\infty, A \}
\]
on trouve
\[
A_i^\circ = A_i - A_{i-1} , \qquad A_\infty^\circ = N , \qquad A^\circ = A
\]
\note{la page porte $A_i - A_{i-1}$ ; à la p.~52 c'était $A_i \smallsetminus A_{i+1}$ ; la suite étant décroissante, on attendrait $i+1$.}
tous ces ens.\ sont $\neq \emptyset$ et mutuellement disjoints, réunion est
$|\Phi| = A_0 \cup A_\infty = A_0 \amalg A$. On a ici $A_\infty^\circ \subset A_i$
pour $i \neq \circ$\note{sic, sans doute $i \neq \infty$.}, alors que
$A_\infty \not\subset A_i$.

\section{Homomorphismes d'ateliers}
\note{titre de l'auteur, souligné, p.~54 ; au-dessus, entre « Homomorphismes » et « d'ateliers », il insère « de magasins et ».}

\page{54}
\struck{Soient}

\textbf{Ex~1} Soit $f : \mathcal{L}' \to \mathcal{L}$ une application
ensembliste, on va en déduire une application canonique
\[
f^* : \mathrm{Figures}(\mathcal{L}) \longrightarrow \mathrm{Figures}(\mathcal{L}') .
\]
Pour toute figure \add{une $\Phi$} dans $\mathcal{L}$, on pose
\[
\text{(1)} \qquad f^*(\Phi) = \{ f^{-1}(A) \mid A \in \Phi,\ f^{-1}(A^\circ) \neq \emptyset \}
\]
\[
\text{i.e. } A^\circ \cap \mathcal{L}_0 \neq \emptyset, \quad \text{où } \mathcal{L}_0 = f(\mathcal{L}')
\]
On peut considérer $f^*$ comme un composé de \struck{deux} \add{trois} applications ;
l'une
\[
A \longmapsto \text{\struck{$A$}}\ F_{\mathcal{L}_0} \cdot A
\]
\marginal{$F_{\mathcal{L}_0} = \{\mathcal{L}_0\}$, figure réduite à la seule strate $\{\mathcal{L}_0\}$ si $\mathcal{L}_0 \neq \emptyset$ i.e.\ $\mathcal{L}' \neq \emptyset$, sinon $F_{\mathcal{L}_0} = \emptyset$}
va de
\[
\mathrm{Figures}(\mathcal{L}) \longrightarrow \mathrm{Figures}(\mathcal{L})_{\mathcal{L}_0} ,
\]
ce dernier ens.\ désignant le sous-ens.\ des premières formé des figures
\uncertain{ayant support} dans $\mathcal{L}_0$. On a d'autre part
(\ill{})
\[
\mathrm{Figures}_{\mathcal{L}_0}(\mathcal{L}) \xrightarrow{\ \sim\ } \mathrm{Figures}(\mathcal{L}_0)
\]
et enfin
\[
f_0^* : \mathrm{Figures}(\mathcal{L}_0) \longrightarrow \mathrm{Figures}(\mathcal{L}') ,
\]
où cette fois, vu que $f$ est surjectif, on peut poser simplement
\[
\text{(2)} \qquad f_0^*(\Phi_0) = \{ f^{-1}(A) \mid A \in \Phi_0 \} ,
\]
la condition $f_0^{-1}(A^\circ) \neq \emptyset$ étant automatique
\note{la phrase se poursuit p.~55.}
\marginal{\ill{} l'ens.\ \ill{} cas \ill{} cas (\ill{}) \ill{} : $\mathrm{Figures}(\mathcal{L}) \to \mathrm{Figures}$ \ill{} \ill{} \ill{} \ill{} \ill{} \ldots}
\marginal{$f_0$ \ill{} l'appl.\ (\ill{}) \ill{} $\mathcal{L}' \to \mathcal{L}_0$}

\page{55}
pour tout $A \in \Phi_0$.

L'application (\ill{}) peut être vue comme un «~passage au quotient~», obtenu
en posant dans $\mathrm{Fig}(\mathcal{L})$ formellement «~$X = \emptyset$~»
pour toute multistructure $X$ telle que
\[
X^\circ \cap \mathcal{L}_0 = \emptyset \quad \text{i.e.} \quad F_{\mathcal{L}_0} \mathbin{|{\circ}|} X
\]
\struck{Il y a deux \ill{},} \struck{une} (\uncertain{en \ill{} attention}, on
conserve les sous-strates $Y$ pour lesquelles $Y^\circ \cap \mathcal{L}_0 \neq \emptyset$).
\struck{Il y} On devine que de façon générale, un tel passage au quotient
\uncertain{revient} \uncertain{à se donner}, dans un magasin $M$, une partie qui
soit $|{\circ}|$-fermée, ou un «~$|{\circ}|$~»-support, dont les éléments seront
«~négligés~» : $\circ$. Si on néglige $X$, on ne néglige pas pour autant toutes ses
strates, mais du moins toutes les $Y$ telles que $Y \mathring{\ll} X$
(\uncertain{moralement}, telles que $Y^\circ \subset X^\circ$)
\marginal{\textbf{NB}}

\struck{L'application}
L'application (\ill{}) est un isomorphisme, inverse de
\[
\mathrm{Figures}_{\mathcal{L}_0}(\mathcal{L}_0) \xrightarrow{\ \sim\ } \mathrm{Figures}(\mathcal{L})_{\mathcal{L}_0} \hookrightarrow \mathrm{Figures}(\mathcal{L}) ,
\]
et le composé
\[
i_* : \mathrm{Figures}(\mathcal{L}_0) \longrightarrow \mathrm{Figures}(\mathcal{L})
\]
est \struck{\ill{}} «~plongement~», d'image $\mathrm{Figures}(\mathcal{L})_{\mathcal{L}_0}$.

\page{56}
Il a l'air d'avoir les propriétés similaires de celles de $f_0^*$ (qui est elle
aussi un plongement) avec une différence importante. Dans les ateliers
considérés, il y a un plus \uncertain{petit} élément $F_{\mathcal{L}}$,
$F_{\mathcal{L}'}$, $F_{\mathcal{L}_0}$. Dans les cas
\add{$f^*$ (1), $f_0^*$ (\uncertain{ou} $i_0^*$)} (\ill{}) et (\ill{}),
\struck{\ill{}} (cas d'une application «~surjective~») l'élément unité est
transformé en itou, mais pas dans le cas d'une application «~non
surjective~» par une inclusion. Du pt de vue supports, on devine que la
différence cruciale, c'est que dans le cas $\varphi^*$, le support
\struck{\ill{}} total est transformé en itou, mais pas dans le cas
$\varphi_*$. Et encore : $\varphi^*$ commute (sur les supports) au passage
au complémentaire, et $\varphi_*$ pas.

Notons que dans le cas d'une inclusion $\varphi = i : \mathcal{L}_0 \hookrightarrow \mathcal{L}$,
\add{$i$} \uncertain{on a vu} à la fois avec $i^*$, et avec $i_*$, on a
\[
i^* i_* = \mathrm{id}_{\mathcal{F}_0} , \qquad
\text{\struck{$i_* i^*$}}\ i_* i^*(X) = i_*(\mathbb{1}) \cdot X , \qquad \text{plus généralement}
\]
\[
i_*(X') \cdot Y = i_*(X' \cdot i^*(Y))
\qquad \text{(formule de projection)}
\]
\note{le facteur $i_*(\mathbb{1})$ est lu sur une surcharge.}
À la fois les $\varphi^*$ et les $\varphi_*$ commutent aux $\mathrm{Sup}^s$
quelconques (y inclus $\varphi^*(\emptyset) = \emptyset$) et $\varphi_*$ aux
$\mathrm{Inf}^s$ de familles non vides.

\page{57}
Il y a une difficulté avec $f^*$, i.e.\ (qui ne se présente pas pour $f_0^*$,
$i_*$) : il est bien vrai que si $X$, $Y$ sont des multistructures, $X'$, $Y'$
leurs images, on a
\[
X \mathring{\ll} Y \Longrightarrow X' \mathring{\ll} Y'
\]
mais pas
\[
X \leq Y \Longrightarrow X' \leq Y' \quad \text{ou seulement } X' \ll Y'
\]
puisqu'il peut arriver que $Y' = \emptyset$, $X' \neq \emptyset$ ! Cela montre donc
que les \struck{images} \add{applications} «~inverses~» $f^*$,
\struck{celles qui ne sont pas des «~plongements~»} pour $f$ non surjectif,
\struck{\ill{} doivent être traitées avec} mais qui impliquent un peu des
passages au quotient, doivent être traitées avec \add{\uncertain{soin}}
\struck{\uncertain{attention}}. Elles ne sont pas justiciables d'une notion
pour-bonhomme de «~morphismes~» de magasins ou d'ateliers, où on aurait
\[
X \mathring{\ll} Y \Longrightarrow X' \mathring{\ll} Y' \quad \text{mais aussi}
\]
\[
X \leq Y \Longrightarrow X' \leq Y' \quad \text{donc aussi}
\]
\[
X \ll Y \Longrightarrow X' \ll Y'
\]
Mais dans tous les cas, \uncertain{morphismes} ou \uncertain{non}, \uncertain{pourtant}
\[
X \mathbin{|{\circ}|} Y \Longrightarrow X' \mathbin{|{\circ}|} Y' .
\]
\note{« pour-bonhomme » est lu tel quel ; le mot est incertain.}

\page{58}
Le point délicat ici, est apparemment dans la notion de passage au quotient,
\struck{plus \ill{} \ill{} contraction} qui nous fait sortir de la notion
d'homomorphismes pousse-bouton.

Dans tous ces exemples, une multistructure \uncertain{était} transformée en
itou (cas de \add{$f^*$ pour} $f$ surjectif, ou de $i_*$), ou bien en itou ou
$\emptyset$ (cas de $f^*$ en général). Il y a des cas qui s'apparentent aux
«~plongements~» d'un atelier dans un autre, où multistructure n'est
\struck{\ill{}} pas transformée en itou :

\textbf{Exemple~2} Soient $\mathcal{L}$, $\mathcal{L}'$ deux ens.,
$\mathcal{L}' \neq \emptyset$ considérons $F' \in \mathrm{Figures}(\mathcal{L}')$,
\add{$F' \neq \emptyset$,} et
\[
\mathrm{Fig}(\mathcal{L}) \longrightarrow \mathrm{Fig}(\mathcal{L} \times \mathcal{L}')
\]
\[
F \longmapsto F \times F' = \mathrm{pr}_1^*(F) \cdot \mathrm{pr}_2^*(F')
\]
C'est le composé d'un $\mathrm{pr}_1^*$ (par déf.\ régulier, $\mathrm{pr}_1$
surjectif), et d'une application interne à $\mathrm{Fig}(\mathcal{L} \times \mathcal{L}')$,
savoir la
\[
\Phi \longmapsto \Phi \cdot F \qquad (\text{où } F = \mathrm{pr}_2^*(F')) .
\]
C'est là un autre exemple d'une application qui (si $\mathrm{supp}\,F$ est
\uncertain{tout} \ill{}) a \ill{} \ill{} \ill{} \ill{} les \ill{} multistructures
\ill{} \ill{}
\note{les deux dernières lignes, serrées au bas de la feuille avec des ajouts
interlinéaires, ne se lisent que par fragments.}

\page{59}
\textbf{Exemple~3} Passage d'un atelier à un sous-atelier, ou d'un magasin
\add{$M$} à un sous-magasin \add{$M'$}, dans l'esprit de restreindre le type de
multistructures envisagées (par des propriétés de type connexité ou
irréductibilité, régularité, être des «~boules~» ou des «~simplexes~» [dans un
sens plus ou moins fort]) mais qui restent assez nombreuses pour
«~\add{pouvoir} \uncertain{résumer} $M$~» — notamment une des \struck{\ill{}}
\add{$M'$} subdivisions \uncertain{arbitraires}.

Dans un tel cas, \uncertain{sûrement} on s'attendra que pour le plongement
$M' \to M$, toutes les structures
\[
\leq ,\quad \mathring{\ll} ,\quad \mathbin{|{\circ}|} \quad /\!/ \quad \text{\struck{$\ll$}} ,\quad \|
\]
\note{la ligne porte, après $\mathbin{|{\circ}|}$, deux traits obliques puis un $\ll$ barré et deux traits verticaux.}
soient conservées, à l'exception \uncertain{toutefois} de $\mathring{\ll}$ qui serait
remplacé par $\mathring{\ll}_{\mathrm{pol}}$, et par conséquent $\ll$ par
$\ll_{\mathrm{pol}}$ \struck{$\ll_{\mathrm{rel}}$}.
\note{l'indice « pol » est lu deux fois ; sous le dernier $\ll$, biffé, un indice que nous lisons « rel ».}
Donc on est ici dans le cas d'une notion d'homomorphisme \emph{de magasins},
qui est un plongement — donc pas de difficultés du type de celles évoquées avec
les exemples~1.

\page{60}
Cela nous conduit au passage avec la question de la stabilité des
\add{(des magasins)} \uncertain{axiomes}, par passage : à un sous-magasin
(d'\uncertain{axiomatique}, pour un atelier \uncertain{associé} au magasin, un
sous-atelier correspondant).

Soit donc $(M, \leq, \mathring{\ll}, \mathbin{|{\circ}|})$ un magasin, $M'$
un sous-ens., que je munis des relations induites. Je suppose
\begin{itemize}
  \item[a)] $M'$ \uncertain{fermé} dans $M$, $\leq$ \uncertain{i.e.} le
  \uncertain{membre} des \uncertain{dessus} ! $\mathrm{Mag}\,1^*$ (trivial,
  p.~22), $\mathrm{Mag}\,2^*$ a) b) aussi\add{, grâce à a)}. Mais on voudrait
  que $\ll_{M'}$ soit $\ll_M \mid M'$, ce qui amène à poser
  \item[b)] Si $X \mathring{\ll} Y' \leq Y$, alors $X, Y \in M' \Rightarrow Y' \in M'$
\end{itemize}
$\mathrm{Mag}\,3$ trivial, $\mathrm{Mag}\,4$ aussi. Notons aussi que
\[
\|_{M'} = \|_M \mid M' , \qquad \ll_{M'} = \ll_M \mid M .
\]
Tout marche !
\note{« p.~22 » renvoie à la pagination de l'auteur, hors de ce lot. Dans la
dernière formule le second $\ll$ porte une barre, et la restriction est écrite
$\mid M$ (sic) ; on attendrait $\mid M'$. L'étude du sous-magasin se poursuit
après la fin du lot.}

\end{document}
